% Part: many-valued-logic
% Chapter: infinite-valued-logics
% Section: introduction

\documentclass[../../../include/open-logic-section]{subfiles}

\begin{document}

\olfileid{mvl}{inf}{int}

\olsection{مقدمة}

لا يلزم أن يكون عدد قيم الصدق في المصفوفة منتهيًا. ومن الخيارات
الواضحة لمجموعة غير منتهية من قيم الصدق مجموعة الأعداد النسبية بين
$0$ و$1$، أي~$V_\infty = [0,1] \cap \Rat$؛ وبعبارة أخرى،
\begin{align*}
    V_\infty & = \Setabs{\frac{n}{m}}{n,m \in \Nat \text{ و } m>0 \text{ و } n\le m}.
\intertext{وعند النظر في مجموعة قيم الصدق غير المنتهية هذه، يفيد غالبًا
النظر أيضًا في المجموعات الجزئية}
V_m & = \Setabs{\frac{n}{m-1}}{n \in \Nat \text{ و } n<m}
\intertext{فمثلًا، $V_5$ هي المجموعة ذات قيم الصدق الخمس المتساوية التباعد،}
V_5 & = \{0, \frac{1}{4}, \frac{1}{2}, \frac{3}{4}, 1\}.
\end{align*}
في المنطقات القائمة على مجموعات قيم الصدق هذه، لا تكون عادة إلا~$1$
مميزة؛ أي~$V^+ = \{1\}$. وبعبارة أخرى، نجعل~$1$ تؤدي دور الصدق
(المطلق)، و$0$ دور الكذب المطلق، لكن~!!{formula}s قد تأخذ أي قيمة
متوسطة في~$V$.

ويمكن أيضًا النظر في المجموعة~$V_{[0,1]} = [0,1]$ لجميع الأعداد
\emph{الحقيقية} بين~$0$ و$1$، أو في مجموعات جزئية غير منتهية أخرى من
$[0,1]$. وغالبًا ما تسمى المنطقات ذات مجموعة قيم الصدق هذه
\emph{ضبابية}.


\end{document}

